% !TeX root = ../组合数学笔记.tex
% 由主文档通过 \input 载入；请编译主文档。
\setcounter{chapter}{3} % 保留原书章号。
\chapter{抽屉原理与存在性证明}
\label{ch:pigeonhole}

计数也能用于证明某种结构必然存在。如果对象太多，而在不出现目标结构时每类只能容纳很少对象，就会产生矛盾。这里的困难通常不是求和，而是设计恰当的分类。

\section{基本形式与容量形式}

\begin{theorem}[抽屉原理（pigeonhole principle）]
将 $N\ge0$ 个对象放入 $k\ge1$ 个盒子，至少一个盒子含有不少于 $\lceil N/k\rceil$ 个对象。更一般地，若第 $i$ 盒的容量为非负整数 $c_i$，且 $N>\sum_{i=1}^k c_i$，则至少一盒超过其容量。
\end{theorem}
\begin{proof}
若每盒均不超过容量，总数至多为容量之和，得一般形式。对基本形式，若每盒都少于 $\lceil N/k\rceil$，每盒对象数都严格小于 $N/k$，总数便小于 $N$，矛盾。
\end{proof}

写 $N=qk+r$，其中 $0\le r<k$。使 $r$ 个盒子各有 $q+1$ 个对象，其余各有 $q$ 个，就达到上述下界。因此在没有附加结构时，不能普遍保证更大的占用数。证明阈值最优通常需要两部分：超过阈值必然发生，以及在阈值以下确实存在避免目标的构造。

\section{余数、区间与部分和}

两个整数除以 $m\ge1$ 的余数相同，当且仅当其差被 $m$ 整除。因此从 $m+1$ 个整数中必能找到两个差被 $m$ 整除。若问题涉及连续区间的和，应先把区间和改写为部分和之差。

\begin{proposition}
对任意 $n\ge1$ 个整数 $a_1,\ldots,a_n$，存在 $1\le i\le j\le n$，使 $a_i+\cdots+a_j$ 被 $n$ 整除。
\end{proposition}
\begin{proof}
置 $s_0=0$、$s_j=\sum_{t=1}^j a_t$。$n+1$ 个部分和只有 $n$ 种可能余数，故有 $0\le u<v\le n$ 满足 $s_u\equiv s_v\pmod n$。于是非空区间之和 $a_{u+1}+\cdots+a_v=s_v-s_u$ 被 $n$ 整除。
\end{proof}

类似地，把 $[0,1]$ 划成 $m$ 个长度为 $1/m$ 的区间，采用左闭右开、最后一段包含右端点的约定，则任意 $m+1$ 个点中至少两个落在同一段，距离不超过 $1/m$。端点约定保证每个点只属于一个盒子；它是证明的一部分。

\section{Erd\H{o}s--Szekeres 单调子序列定理}

\begin{theorem}
设 $r,s\ge2$。任意由 $(r-1)(s-1)+1$ 个互不相同实数组成的序列，含有长度为 $r$ 的严格递增子序列或长度为 $s$ 的严格递减子序列。
\end{theorem}
\begin{proof}
对位置 $i$，令 $p_i,q_i$ 分别为以该位置结尾的最长严格递增、严格递减子序列的长度。若结论不成立，则所有 $(p_i,q_i)$ 落在 $[r-1]\times[s-1]$ 中。对 $i<j$，若 $a_i<a_j$，把 $a_j$ 接到以 $i$ 结尾的递增子序列后，得 $p_j\ge p_i+1$；若 $a_i>a_j$，同理 $q_j\ge q_i+1$。故各有序对互不相同，但位置数超过可用有序对数，矛盾。
\end{proof}

子序列只要求下标递增，不要求相邻。阈值不能降低：把 $1,\ldots,(r-1)(s-1)$ 分成 $r-1$ 个连续数值块，每块含 $s-1$ 个数；块按数值从小到大排列，块内递减。递增子序列每块至多取一个，递减子序列只能在一块内，因而两种目标均不存在。

\section{极端原理、不变量与构造}

\emph{极端原理（extremal principle）}通过选择某个量最大的、最小的或最短的对象，利用其不可进一步改进的性质。它要求极端对象确实存在：有限非空集合自动满足这一点，一般无限集合则未必如此。\emph{不变量（invariant）}是在每次合法操作后保持不变的量；若初态与目标态的不变量不同，就证明目标不可达。

\begin{example}[棋盘染色]
从 $8\times8$ 棋盘删除两个对角顶角后，不能用 $1\times2$ 骨牌覆盖。将棋盘交替染成黑白两色，每块骨牌必覆盖一黑一白；两个对角顶角同色，删除后两种颜色的格数分别为 $30$ 与 $32$，而任何骨牌覆盖要求二者相等，故不可能。这里证明的是不存在，并未依靠逐一枚举铺法。
\end{example}

\begin{example}[同一集合的两种抽屉]
从 $[2n]$ 中选出 $n+1$ 个数。按 $\{1,2\},\{3,4\},\ldots,\{2n-1,2n\}$ 分组，某组取了两个数，故有相邻数。若要证明一个整除另一个，改为按奇数部分分组：每个正整数唯一写成 $2^a u$，其中 $u$ 为奇数；$[2n]$ 中只有 $n$ 种可能的奇数部分。同组的两个数为 $2^a u,2^b u$，较小者整除较大者。不同目标要求不同抽屉。
\end{example}

\section{习题}
\begin{enumerate}
  \item\label{ex:pigeonhole-1} 证明任意六个整数中有两个之差被 $5$ 整除，并构造五个整数使任意两数之差都不被 $5$ 整除。
  \item\label{ex:pigeonhole-2} 在单位正方形内取五个点，证明有两点距离不超过 $\sqrt2/2$；说明边界点怎样归类。
  \item\label{ex:pigeonhole-3} 构造长度为 $9$ 的不同整数序列，使其没有长度为 $4$ 的递增或递减子序列。
  \item\label{ex:pigeonhole-4} 初始写有整数 $1,2,\ldots,10$，每次擦掉两个数 $a,b$ 并写入 $a+b$。证明最后一个数不依赖操作次序；若改为 $a+b+1$，最后结果是多少？
\end{enumerate}

